%% Section 4: The Yang-Baxter Groupoid and the Five-Vertex Classification
%% To be \input{} into the main document.

\section{The Yang-Baxter Groupoid and the Five-Vertex Classification}
\label{sec:groupoid}

Section~\ref{sec:lattice} established that the KTW puzzle model is an
integrable vertex model: the puzzle pieces carry R-matrix weights
satisfying the Yang-Baxter equation, and the LR coefficient is a
partition function.  We also saw that the Schur case is not isolated
--- Grothendieck, Hall-Littlewood, and Demazure characters each have
their own solvable (or quasi-solvable) lattice model, each with its
own R-matrix.  The section closed with two questions: what is the
\emph{relationship} between the many combinatorial models for the same
coefficient, and what is the \emph{relationship} between the R-matrices
at different levels of the character hierarchy?

This section answers both questions.  The first answer is the
\emph{Yang-Baxter groupoid} of Bump and
Naprienko~\cite{bump-naprienko-2025}: classical bijections between
combinatorial models for LR coefficients are morphisms in a groupoid
whose composition law is the Yang-Baxter equation.  The second answer
is a \emph{classification} of five-vertex R-matrices satisfying the
YBE, which organises the R-matrices into a strict three-level
hierarchy.  Together, these two results set the stage for the
categorical phase diagram of Section~\ref{sec:hierarchy}.


\subsection{The groupoid framework of Bump and Naprienko}
\label{sec:groupoid-bn}

The starting point is a simple observation: if the Yang-Baxter
equation~\eqref{eq:ybe} holds, then replacing the crossing pattern
$R_{12}R_{13}R_{23}$ with $R_{23}R_{13}R_{12}$ in a portion of the
lattice does not change the partition function.  Each such replacement
is a \emph{local bijection} on configurations --- it permutes the
internal spin assignments without altering the boundary data or the
total weight.  Bump and Naprienko~\cite{bump-naprienko-2025} elevate
this observation to a categorical framework: the \emph{Yang-Baxter
groupoid}.

\subsubsection*{Objects and morphisms}

The objects of the Yang-Baxter groupoid are \emph{models}: each model
is specified by a choice of lattice domain (the shape of the grid),
boundary data (the partitions or binary strings fixed on the boundary
edges), and local vertex weights (the R-matrix entries at each
vertex).  Two models may have different lattice domains --- a
triangular puzzle grid versus a rectangular grid of Gelfand-Tsetlin
type, say --- but both compute the same structure constants through
their respective partition functions.

The morphisms are \emph{sequences of Yang-Baxter moves}.  A single YBE
move acts on a local patch of the lattice containing three crossing
lines: it replaces the crossing order $(1,2), (1,3), (2,3)$ with
$(2,3), (1,3), (1,2)$, or vice versa.  Because the YBE guarantees
that both orderings produce the same map on boundary spins, the move
defines a bijection between the configurations of the original local
patch and the configurations of the transformed patch, weighted
consistently.  A \emph{groupoid morphism} is a finite composition of
such local moves, carrying one model to another while preserving the
partition function at every step.

The groupoid axioms are satisfied by construction:
\begin{itemize}
\item \emph{Identity.}  The trivial sequence of zero moves is the
  identity morphism on any model.
\item \emph{Composition.}  Composing two sequences of YBE moves
  gives another sequence.
\item \emph{Inverses.}  Every YBE move is its own inverse (applying
  it twice returns to the original crossing order), so every morphism
  is invertible.
\end{itemize}

\noindent
The crucial additional property is \emph{coherence}: if two different
sequences of YBE moves relate the same pair of models, they induce the
\emph{same} bijection on configurations.  This follows from the
Zamolodchikov tetrahedron equation in low rank and from braid group
relations in general --- the braid group acts faithfully on the space
of crossing orders, and the YBE is precisely the braid relation.  In
practice, coherence means that groupoid morphisms are well-defined
maps, not merely equivalence classes of move sequences.

\begin{remark}
The groupoid is not a group: different objects (models) are not
isomorphic in general --- a puzzle model on a triangular lattice and
a Gelfand-Tsetlin model on a rectangular lattice have genuinely
different configuration spaces.  But whenever two models \emph{are}
connected by a morphism, the morphism provides a canonical,
weight-preserving bijection between their configuration sets.
\end{remark}

\subsubsection*{Recovering classical bijections}

The power of the groupoid framework lies in its ability to
\emph{explain} bijections that were previously understood only through
explicit, often intricate constructions.  Bump and Naprienko show that
several classical bijections in the theory of LR coefficients are
Yang-Baxter groupoid morphisms:

\begin{enumerate}
\item \textbf{The commutation isomorphism.}  The symmetry
  $c^{\nu}_{\lambda,\mu} = c^{\nu}_{\mu,\lambda}$ is obvious from the
  Schur function definition ($s_\mu \cdot s_\nu = s_\nu \cdot s_\mu$
  by commutativity), but at the level of puzzle tilings it is
  nontrivial: the two sides correspond to different boundary
  conditions, and one must construct an explicit bijection between the
  two sets of tilings.  In the groupoid framework, this bijection is a
  morphism obtained by ``sliding'' one set of auxiliary lines past
  another using YBE moves.  The commutativity of multiplication
  becomes a geometric consequence of integrability.

\item \textbf{The Sch\"utzenberger involution.}
  Gustafsson and Westerlund~\cite{gustafsson-westerlund-2025} showed
  that the Sch\"utzenberger involution on semistandard Young
  tableaux --- an involution that plays a central role in the theory of
  jeu de taquin and crystal bases --- arises as a Yang-Baxter groupoid
  morphism.  Their construction interprets the involution as a sequence
  of R-matrix moves that ``reverses the spectral parameters,''
  transforming the lattice model by reflecting the order in which rows
  are processed.  The fact that this reflection preserves the
  partition function is a direct consequence of transfer matrix
  commutativity~\eqref{eq:transfer-commute}.

\item \textbf{The puzzle--tableau correspondence.}  The bijection
  between KTW puzzle tilings and LR skew-tableaux, mediated by
  Purbhoo's mosaics~\cite{purbhoo-2007}, has a natural interpretation
  as a groupoid morphism.  The two models --- puzzles on a triangular
  domain and tableaux on a skew shape --- correspond to different
  objects in the groupoid, and the bijection between them is
  a composition of YBE moves that deforms the triangular lattice into
  the rectangular lattice underlying the tableau model.  We state this
  as a guiding principle rather than a proven theorem: verifying it
  requires checking that each step of Purbhoo's mosaic
  correspondence can be decomposed into elementary R-matrix moves, a
  computation that has been carried out for small examples ($n \leq 4$)
  but not yet in full generality.
\end{enumerate}

The unifying message is this: the ``mysterious'' bijections of LR
coefficient theory are not feats of combinatorial ingenuity.  They are
canonical consequences of the Yang-Baxter equation --- different
paths through the groupoid, all connecting the same pair of models,
all inducing the same bijection by coherence.

\subsubsection*{Connection to Gaetz-Gao}

The scope of the groupoid extends beyond the classical models.
Gaetz and Gao~\cite{gaetz-gao-2024} proved that every
\emph{integrable triangle-accessible} (ITA) lattice model --- a class
they define by axiomatising the properties needed for a vertex model to
compute LR-type structure constants --- bijects canonically with a KTW
puzzle model.  Their result implies that puzzles are, in a precise
sense, a \emph{universal} object in the groupoid: every ITA model is
connected to the puzzle model by a groupoid morphism.

This universality is satisfying but also raises a question.  If the
puzzle model is canonical within the groupoid, what determines which
R-matrices \emph{define} the groupoid in the first place?  Which
$4 \times 4$ matrices satisfy the Yang-Baxter equation?  For the
five-vertex models relevant to LR coefficient theory, this is a finite
classification problem.  We turn to it now.


\subsection{The five-vertex classification and chirality}
\label{sec:groupoid-classification}

The R-matrices of Section~\ref{sec:lattice-rmatrix} are $4 \times 4$
matrices acting on $\mathbb{C}^2 \otimes \mathbb{C}^2$, with edge
spins in $\{0, 1\}$.  The conservation law (total input spin equals
total output spin) restricts the nonzero entries to the six
``ice-type'' configurations:
\[
  \begin{array}{cccc}
    (0,0) \to (0,0), &
    (1,1) \to (1,1), &
    (0,1) \to (0,1), \\[4pt]
    (1,0) \to (1,0), &
    (0,1) \to (1,0), &
    (1,0) \to (0,1).
  \end{array}
\]
A \emph{six-vertex model} assigns nonzero weights to all six; a
\emph{five-vertex model} forbids exactly one of the six, setting its
weight to zero.  The KTW puzzle model is a five-vertex model: it
forbids the configuration $(1,0) \to (1,0)$, in which a spin-$1$
entering from the left and a spin-$0$ entering from below pass through
the vertex unchanged.

\subsubsection*{The classification question}

Among all five-vertex R-matrices with a spectral parameter $u$, which
ones satisfy the Yang-Baxter equation~\eqref{eq:ybe}?  Since there are
six choices of which vertex to forbid, and for each choice the
remaining five weights are functions of $u$ to be determined, this is
a finite (but nontrivial) classification problem.

The YBE is a system of 64 cubic equations in the vertex weights.
Most of the 64 equations are trivially satisfied when one of the six
vertices is forbidden.  The content of the classification is in the
remaining equations: they impose algebraic relations among the five
nonzero weights, and these relations either have nontrivial solutions
(the model is integrable) or they do not (the model is not
integrable).

\subsubsection*{The classification result}

The outcome of the classification is as follows.

\begin{enumerate}
\item For each choice of forbidden vertex, the five-vertex YBE system
  admits a family of solutions parameterised by the spectral parameter
  $u$ and a finite number of structural constants (the values of the
  weights at a reference point).

\item Up to the symmetries of the problem (transposition of input and
  output, reversal of spin labels, rescaling of the R-matrix), the
  solutions organise into families labelled by which vertex is
  forbidden.  The puzzle model (forbidding $(1,0) \to (1,0)$) is one
  such family.

\item There is essentially one nontrivial five-vertex solution
  satisfying the full spectral-parameter YBE:
  \begin{equation}
  \label{eq:five-vertex-unique}
    R(u) \;=\; c(u)\bigl(P + \alpha u \cdot E_{22}\bigr),
  \end{equation}
  where $P$ is the permutation (swap) matrix, $E_{22}$ is the matrix
  unit with a $1$ in the $(2,2)$ position and zeros elsewhere,
  $\alpha$ is a structural constant, and $c(u)$ is a scalar
  normalisation.  This R-matrix computes $\omega$-dual Jacobi-Trudi
  determinants --- the $t = 0$ specialisation of the
  Wheeler-Zinn-Justin Hall-Littlewood
  model~\cite{wheeler-zinn-justin-2016}.
\end{enumerate}

\subsubsection*{Chirality}

To state the classification precisely, we need the concept of
\emph{chirality}.  Define the two ``non-crossing'' vertices as those
where each spin passes straight through: $(0,1) \to (0,1)$ with
weight $a_1$, and $(1,0) \to (1,0)$ with weight $a_2$.  The two
``crossing'' vertices are $(0,1) \to (1,0)$ with weight $b_1$ and
$(1,0) \to (0,1)$ with weight $b_2$.  A vertex model is
\emph{chiral} if $b_1 \neq b_2$: the weights for left-to-right and
right-to-left crossings differ.  It is \emph{label-asymmetric} if
$a_1 \neq a_2$: the two non-crossing vertices have different weights.

In a five-vertex model, one of the six weights is set to zero.
Forbidding one of the crossing vertices (say $b_2 = 0$, as in the KTW
puzzle model) automatically introduces chirality: the two crossing
types are no longer equivalent, since one is allowed and the other is
forbidden.  This raises the question: does chirality obstruct the
Yang-Baxter equation?

The answer, perhaps surprisingly, is \emph{no}.  The R-matrix
\eqref{eq:five-vertex-unique} is chiral (one crossing type is
suppressed relative to the other, as the permutation matrix $P$ treats
the two asymmetrically once $E_{22}$ is added) and it satisfies the
full YBE.  The KTW puzzle model itself is chiral and integrable.
Chirality alone is not an obstruction to integrability.

What \emph{is} an obstruction is the \emph{combination} of chirality
with label asymmetry.  The complete five-vertex classification can be
stated as a dichotomy:

\begin{quote}
\textbf{Chirality principle.}  A five-vertex model with spectral
parameter satisfies the Yang-Baxter equation if and only if either
\begin{enumerate}
\item[(a)] the non-crossing weights are equal: $a_1 = a_2$; or
\item[(b)] the model is a degenerate (free-fermionic) specialisation
  of the six-vertex model.
\end{enumerate}
In particular, chirality combined with $a_1 \neq a_2$ is the
\emph{only} obstruction to the five-vertex YBE.
\end{quote}

This principle was verified computationally for all six choices of
forbidden vertex: of the 64 matrix-element equations comprising the
YBE, exactly 4 produce nonzero residuals in the generic $a_1 \neq a_2$
case, and the residuals are proportional to $q(p - q)$, where $p$ and
$q$ are the two non-crossing weights.  The obstruction vanishes if and
only if $p = q$ (condition (a)) or the model degenerates further
(condition (b)).

The chirality principle has a clear physical interpretation.  In a
chiral model, paths through the lattice ``know'' whether they are
turning left or right --- the model distinguishes orientation.  This
is compatible with the YBE as long as the ``straight-through'' weights
treat both spin values equally ($a_1 = a_2$).  But if the straight
paths \emph{also} distinguish between spin values ($a_1 \neq a_2$),
then the combination of orientational asymmetry (chirality) and spin
asymmetry (label asymmetry) produces an overdetermined system with no
solutions.

\begin{remark}
The chirality principle explains why the KTW puzzle model is
integrable: although it is chiral (the forbidden vertex
$(1,0) \to (1,0)$ breaks the crossing symmetry), its non-crossing
weights are both equal to~$1$, so $a_1 = a_2$.  Condition (a) is
satisfied.  The K-theoretic deformation of
Section~\ref{sec:lattice-extensions} preserves this equality (both
non-crossing weights deform together), which is why the Grothendieck
R-matrix also satisfies the YBE.
\end{remark}


\subsection{The three-level hierarchy}
\label{sec:groupoid-hierarchy}

The five-vertex classification reveals that the space of integrable
models is not a continuum but a stratified landscape.  The R-matrices
satisfying the full YBE form one stratum; below them sit R-matrices
that satisfy a weaker form of the equation; below those, the
permutation matrix $P$, which satisfies the YBE trivially.  This
stratification organises into a strict three-level hierarchy that
connects the classification of this section to the categorical phase
diagram of Section~\ref{sec:hierarchy}.

\subsubsection*{Level 1: Solvable (full YBE)}

At the top of the hierarchy sit the \emph{solvable} models: those
whose R-matrix $R(u)$ satisfies the full spectral-parameter
Yang-Baxter equation~\eqref{eq:ybe}.  The transfer matrices $T(u)$
form a commuting family, the model can be diagonalised by the Bethe
ansatz, and R-matrix moves provide local bijections between
configurations.  The Yang-Baxter groupoid of
Section~\ref{sec:groupoid-bn} lives at this level.

The characters computed by solvable models are the classical symmetric
functions and their deformations: Schur polynomials (rational
five-vertex model), Grothendieck polynomials (trigonometric
five-vertex), and Hall-Littlewood polynomials (higher-spin models).
In each case, the structure constants are nonnegative --- a property
that can be traced to the existence of a partition function with
positive Boltzmann weights.

\subsubsection*{Level 2: Quasi-solvable (factored YBE)}

The middle level contains the \emph{quasi-solvable} models, a term
due to Buciumas and Scrimshaw~\cite{buciumas-scrimshaw}.  Here the
full YBE fails at the vertex level: one cannot freely interchange
three crossing lines by local moves.  But a weaker form of
integrability persists.  The row transfer matrices and the column
transfer matrices each satisfy the YBE \emph{independently} --- the
R-matrix factors into separate ``row'' and ``column'' components, each
integrable on its own, but there is no single unified R-matrix
governing the full model.

The characters at this level are the \emph{key polynomials} (Demazure
characters).  Algebraically, the quasi-solvable level corresponds to
the specialisation $q = 0$ of the Hecke algebra parameter, where the
Hecke relation $(T_i - q)(T_i + 1) = 0$ degenerates to
$T_i(T_i + 1) = 0$ and the Demazure operators
$\pi_i = T_i + 1$ become idempotent: $\pi_i^2 = \pi_i$.  The braid
group symmetry of the solvable level degenerates to the \emph{cactus
group} $J_n$ --- a group generated by interval reversals
$s_{[i,j]}$ satisfying the relations of Henriques and
Kamnitzer~\cite{henriques-kamnitzer-2006} but \emph{not} the braid
relations.

\subsubsection*{Level 3: Combinatorial (trivial YBE)}

At the base of the hierarchy, the R-matrix is simply the permutation
matrix~$P$: the swap map on $\mathbb{C}^2 \otimes \mathbb{C}^2$.
The YBE is satisfied trivially ($P_{12} P_{13} P_{23} =
P_{23} P_{13} P_{12}$ reduces to the symmetric group relation
$s_i s_j s_i = s_j s_i s_j$), but there is no spectral parameter and
no interesting commuting family.  The ``transfer matrices'' are
combinatorial objects --- crystal operators in the sense of
Kashiwara~\cite{kashiwara-1990} --- and the characters are monomials,
the building blocks from which all higher-level characters are
assembled.

\subsubsection*{The hierarchy is strict}

It is natural to ask whether the intermediate level is genuinely
intermediate, or merely a degeneration of the solvable level that
retains all of its structure.  The answer is that the hierarchy is
\emph{strict}: the quasi-solvable level has less structure than the
solvable level, and the difference is not an artifact of the
specialisation $q = 0$.

The sharpest evidence comes from the \emph{Zinn-Justin tetrahedron
equation} (ZTE), a cubic coherence condition on the R-matrix that
generalises the YBE from two dimensions to three.  At the solvable
level, the ZTE holds (it follows from the full parametric YBE by a
standard argument involving auxiliary spaces).  At the quasi-solvable
level, the ZTE \emph{fails}: a direct computation at rank~$3$ shows
that the residual is proportional to $q^2$ and transforms as an
antisymmetric 2-form on the space of spectral parameters.

This failure has a precise interpretation: it measures the
\emph{curvature} of a connection on a vector bundle over the higher
Bruhat order.  At $q = 0$, the curvature is nonzero, meaning that
parallel transport around a closed loop in the space of row orderings
does not return to the identity.  The holonomy of this connection is
the obstruction to promoting column-level integrability to full
vertex-level integrability.  The quasi-solvable level is a flat
$q = 0$ limit in one direction (each column is integrable) but
\emph{curved} in the transverse direction (the columns do not
interleave coherently).

\subsubsection*{A structural summary}

The three levels, together with their algebraic and categorical
signatures, are summarised in the following table:

\begin{center}
\renewcommand{\arraystretch}{1.4}
\begin{tabular}{llll}
\toprule
\textbf{Level} & \textbf{YBE status} & \textbf{Algebra} & \textbf{Characters} \\
\midrule
Solvable & Full spectral-parameter YBE &
  $H_q$ (generic $q$) &
  Schur, Grothendieck, HL \\
Quasi-solvable & Factored (column-level) YBE &
  $H_0$ ($q = 0$ Hecke) &
  Key / Demazure \\
Combinatorial & Trivial ($R = P$) &
  $\mathbb{C}[S_n]$ &
  Monomials \\
\bottomrule
\end{tabular}
\end{center}

\noindent
The rightmost column reveals a striking alignment: the integrability
level of the R-matrix controls which characters the model computes.
Solvable models compute symmetric functions (Schur and their
deformations); quasi-solvable models compute nonsymmetric but
``partially symmetric'' characters (key polynomials); combinatorial
models compute the fully nonsymmetric monomials.  The level of
integrability is the level of symmetry.

\subsubsection*{Preview of Section~\ref{sec:hierarchy}}

The three levels of the hierarchy correspond to three types of
monoidal categories:
\begin{enumerate}
\item \textbf{Braided monoidal} (solvable): the R-matrix provides a
  braiding --- a natural isomorphism $V \otimes W \xrightarrow{\sim}
  W \otimes V$ satisfying the braid relation.  The symmetry group is
  the braid group $B_n$.
\item \textbf{Coboundary monoidal} (quasi-solvable): the braiding
  degenerates to a \emph{commutor} --- a natural isomorphism that
  satisfies the weaker cactus group relations but not the braid
  relation.  The symmetry group is the cactus group $J_n$.
\item \textbf{Symmetric monoidal} (combinatorial): the commutor is
  the trivial swap $\tau \colon V \otimes W \to W \otimes V$, and
  $\tau^2 = \mathrm{id}$.  The symmetry group is the symmetric group
  $S_n$.
\end{enumerate}

\noindent
The Hecke algebra $H_q(S_n)$ simultaneously hosts all three
structures: at generic~$q$, the Hecke generators provide a braiding;
at $q = 0$, the braiding dies but the coboundary commutor (constructed
from Drinfeld's unitarised R-matrix) survives; at $q = 1$
(equivalently, passing to the group algebra $\mathbb{C}[S_n]$), only
the symmetric swap remains.  The parameter $q$ governs the transition
between levels, and the Hecke algebra is the \emph{transition algebra}
of the hierarchy.

Making this precise --- identifying what is forgotten at each
transition, and proving that the forgetting is irreversible --- is the
content of Section~\ref{sec:hierarchy}.
